% !TEX program = xelatex
\documentclass[10pt,letterpaper]{article}

\usepackage[letterpaper,left=0.52in,right=0.52in,top=0.42in,bottom=0.42in]{geometry}
\usepackage{fontspec}
\usepackage{microtype}
\usepackage{xcolor}
\usepackage{enumitem}
\usepackage{tabularx}
\usepackage{array}
\usepackage{hyperref}
\usepackage{needspace}
\usepackage[normalem]{ulem}

\setmainfont{texgyretermes-regular.otf}[
  BoldFont=texgyretermes-bold.otf,
  ItalicFont=texgyretermes-italic.otf,
  BoldItalicFont=texgyretermes-bolditalic.otf
]

\definecolor{sectionblue}{HTML}{4F7F9F}
\definecolor{linkblue}{HTML}{124EC7}
\definecolor{rulegray}{HTML}{555555}

\hypersetup{
  colorlinks=true,
  linkcolor=linkblue,
  urlcolor=linkblue,
  pdfauthor={Zheng Jiang},
  pdftitle={Zheng Jiang - Curriculum Vitae},
  pdfsubject={Academic curriculum vitae}
}
\urlstyle{same}

\pagestyle{empty}
\raggedbottom
\setlength{\parindent}{0pt}
\setlength{\parskip}{0pt}
\setlength{\tabcolsep}{0pt}
\setlength{\emergencystretch}{1.5em}
\linespread{1.00}

\setlist[enumerate]{
  leftmargin=0.18in,
  labelsep=0.055in,
  itemsep=3.4pt,
  parsep=0pt,
  topsep=2pt,
  partopsep=0pt
}
\setlist[itemize]{
  leftmargin=0.18in,
  labelsep=0.06in,
  itemsep=1.2pt,
  parsep=0pt,
  topsep=1.4pt,
  partopsep=0pt
}

\newcommand{\self}{\uline{Zheng Jiang}}
\newcommand{\cvsection}[2]{%
  \Needspace{5\baselineskip}%
  \vspace{5.5pt}%
  {\color{sectionblue}\bfseries\fontsize{10.8}{12}\selectfont #1}%
  {\normalfont\color{black}\fontsize{9.4}{10.5}\selectfont\ #2}%
  \par\vspace{1.1pt}%
  {\color{rulegray}\hrule height 0.35pt}%
  \vspace{3.5pt}%
}
\newcommand{\educationentry}[4]{%
  \begin{tabularx}{\textwidth}{@{}>{\raggedright\arraybackslash}X >{\raggedleft\arraybackslash}p{1.42in}@{}}
    \textbf{#1} & \textbf{#2} \\
    \textit{#3} & \textit{#4}
  \end{tabularx}\par
}
\newcommand{\pubitem}[3]{%
  \needspace{4.8\baselineskip}%
  \item \textbf{#1}\par\vspace{-0.4pt}%
  #2\par\vspace{-0.4pt}%
  #3%
}
\newcommand{\awardentry}[3]{%
  \item \begin{tabularx}{\dimexpr\linewidth\relax}[t]{@{}>{\raggedright\arraybackslash}X >{\raggedleft\arraybackslash}p{0.55in}@{}}
    \textbf{#1}, #2 & #3
  \end{tabularx}%
}
\newcommand{\teachingentry}[2]{%
  \item \begin{tabularx}{\dimexpr\linewidth\relax}[t]{@{}>{\raggedright\arraybackslash}X >{\raggedleft\arraybackslash}p{1.33in}@{}}
    \textbf{Teaching Assistant}, #1, BUPT & #2
  \end{tabularx}%
}

\begin{document}
\fontsize{9.8}{11.35}\selectfont

\begin{center}
  {\fontsize{20}{22}\selectfont Zheng Jiang}\\[-0.2pt]
  {\fontsize{9.3}{10.5}\selectfont
  \href{tel:+8615678312406}{Phone: +86 156 7831 2406}
  \enspace--\enspace
  \href{mailto:nicezheng.jiang@gmail.com}{Email: nicezheng.jiang@gmail.com}
  \enspace--\enspace
  \href{https://nicezheng.github.io/}{Website: nicezheng.github.io}\\[-0.4pt]
  \href{https://scholar.google.com/citations?user=b60BQmoAAAAJ}{Google Scholar}
  \enspace--\enspace
  \href{https://github.com/nicezheng}{GitHub}
  \enspace--\enspace
  \href{https://www.linkedin.com/in/zheng-jiang-a55502283/}{LinkedIn}
  \enspace--\enspace
  \href{https://x.com/NicezhengJiang}{X}}\\[-0.2pt]
  School of Artificial Intelligence, BUPT
  \enspace--\enspace Beijing, China
  \enspace--\enspace Ph.D. Candidate in Intelligent Science and Technology
\end{center}
\vspace{-4.5pt}

\cvsection{RESEARCH INTERESTS}{}
My research spans AI in social systems and AI for science. In social systems, I study collaboration decisions in open-source communities and geopolitical bias in how large language models represent human values. In AI for science, I develop domain- and physics-informed methods for geography and ocean forecasting.
\begin{itemize}
  \item \textbf{Human--AI Collaboration in Open-Source Communities:} how developer profiles and activity histories shape collaboration decisions by humans and LLMs, and how better information design can support effective and inclusive collaboration.
  \item \textbf{Geopolitical Bias and Pluralistic AI:} how LLMs represent human values across societies, where geopolitical asymmetries emerge, and how AI can better reflect diverse and sometimes conflicting human values.
  \item \textbf{AI for Geography:} domain-informed and human-centered AI that supports geographic research from visual interpretation and expert interaction to field deployment and continuous learning.
  \item \textbf{AI for Ocean Science:} integrating physical ocean dynamics with data-driven models to improve prediction and enable earlier warning of extreme events.
\end{itemize}

\cvsection{EDUCATION}{}
\educationentry{Beijing University of Posts and Telecommunications (BUPT)}{Beijing, China}{Ph.D. in Intelligent Science and Technology}{Sep. 2022 -- Present}
\begin{itemize}
  \item Advisor: Prof. Wei Wang; Co-supervisor: Prof. Yi Wang
\end{itemize}
\vspace{-1.2pt}
\educationentry{Beijing Forestry University (BJFU)}{Beijing, China}{B.E. in Computer Science and Technology}{Sep. 2018 -- Jun. 2022}

\cvsection{PUBLICATIONS}{(* denotes equal contribution)}
\begin{enumerate}
  \pubitem{Verify the Unverifiable: A Platform for AI in Unverifiable Decision Scenarios}
    {\self, Wei Wang, Gaowei Zhang, Yi Wang.}
    {\textit{12th International Conference on Computational Social Science (IC2S2)}, 2026.}

  \pubitem{Investigating Human and LLMs' Decisions in Unverifiable Environments: A Case Study with GitHub Activity Overview}
    {\self, Wei Wang, Gaowei Zhang, Yang Feng, Yi Wang.}
    {\textit{Findings of the Association for Computational Linguistics (ACL Findings)}, 2026.}

  \pubitem{Towards an Early Warning System for Ocean Heat Extremes Through AI-Ocean Dynamics Synergy}
    {\self, Wei Wang, Gaowei Zhang, Yifei Bao, Zengzhou Hao, Lingyu Xu, Suixiang Shi, Lei Wang, Yi Wang.}
    {\textit{35th International Joint Conference on Artificial Intelligence (IJCAI-ECAI)}, AI and Social Good Track, 2026.}

  \pubitem{How Different LLMs Behave in Unverifiable Decision Scenarios?}
    {\self, Wei Wang, Gaowei Zhang, Yang Feng, Yi Wang.}
    {\textit{48th Annual Meeting of the Cognitive Science Society (CogSci)}, 2026.}

  \pubitem{The Value Atlas of AI: Mapping World Human Values in Large Language Models}
    {Xinyi You\textsuperscript{*}, \self\textsuperscript{*}, Zichan Chen, Gaowei Zhang, Wei Wang, Yang Feng, Libo Liu, Yi Wang.}
    {\textit{Preprint}, 2026.}

  \pubitem{SSTODE: Ocean-Atmosphere Physics-Informed Neural ODEs for Sea Surface Temperature Prediction}
    {\self, Wei Wang, Gaowei Zhang, Yi Wang.}
    {\textit{40th Annual AAAI Conference on Artificial Intelligence (AAAI)}, AI for Social Impact Track, 2026.}

  \pubitem{Maintaining the Heterogeneity in the Organization of Software Engineering Research}
    {Yang Yue, \self, Yi Wang.}
    {\textit{IEEE/ACM International Conference on Software Engineering (ICSE)}, Future of Software Engineering Track, 2026.}

  \pubitem{The Paradox of Being Seen: Signaling Contributors' Unobservable Activities Limits Their Success in Open Collaboration}
    {\self, Wei Wang, Liu Wang, Yang Feng, Min Zhang, Libo Liu, Gaowei Zhang, Yi Wang.}
    {\textit{Preprint}, 2025.}

  \pubitem{Deep Learning for Ocean Parameters Prediction: A Systematic Literature Review}
    {Wei Wang, \self, Anqi Lu, Gaowei Zhang, Yi Wang, Lingyu Xu, Suixiang Shi.}
    {\textit{Preprint}, 2025.}

  \pubitem{Commitment \& Cooperation in Social Dilemmas with Diverse Individual Preferences: An Agent-Based Modeling Approach}
    {\self, Luzhan Yuan, Wei Wang, Gaowei Zhang, Yi Wang.}
    {\textit{PLOS ONE}, 2025.}

  \pubitem{A Big Step Forward? A User-Centric Examination of iOS App Privacy Report and Enhancements}
    {Liu Wang, Dong Wang, Shidong Pan, \self, Haoyu Wang, Yi Wang.}
    {\textit{46th IEEE Symposium on Security and Privacy (S\&P)}, 2025.}

  \pubitem{FenGePad---A Tangible Multi-Prompt Interactive Framework for Deep Dune Segmentation}
    {\self\textsuperscript{*}, Haonan Kang\textsuperscript{*}, Zifeng Wu, Wei Wang, Gaowei Zhang, Eerdun Hasi, Xiaohan Sun, Yi Wang.}
    {\textit{28th European Conference on Artificial Intelligence (ECAI)}, Prestigious Applications of Intelligent Systems (PAIS), 2025.}

  \pubitem{DCV\textsuperscript{2}I: Leveraging Deep Vision Models to Support Geographers' Visual Interpretation in Dune Segmentation}
    {Anqi Lu, Zifeng Wu, \self, Wei Wang, Gaowei Zhang, Eerdun Hasi, Yi Wang.}
    {\textit{AI Magazine}, 2024.}

  \pubitem{DuneSA: A SAM-Based Approach with Domain-Specific Knowledge for Aeolian Dune Segmentation}
    {Anqi Lu, \self, Zifeng Wu, Wei Wang, Gaowei Zhang, Eerdun Hasi, Yi Wang.}
    {\textit{4th ACM International Conference on Information Technology for Social Good (GoodIT)}, 2024.}

  \pubitem{FenGe---An Interactive Framework for Improving the Utility of Deep Dune Segmentation in Geographical Tasks}
    {\self, Anqi Lu, Zifeng Wu, Wei Wang, Gaowei Zhang, Eerdun Hasi, Yi Wang.}
    {\textit{27th European Conference on Artificial Intelligence (ECAI)}, 2024.}

  \pubitem{DCV2I: A Practical Approach for Supporting Geographers' Visual Interpretation in Dune Segmentation with Deep Vision Models}
    {Anqi Lu, Zifeng Wu, \self, Wei Wang, Eerdun Hasi, Yi Wang.}
    {\textit{38th AAAI Conference on Artificial Intelligence (AAAI)}, Innovative Applications of Artificial Intelligence (IAAI), 2024. \textbf{Deployed Application Award}.}

  \pubitem{The Dynamics of Cooperation with Commitment in a Population of Heterogeneous Preferences---An ABM Study}
    {Wei Wang, Luzhan Yuan, \self, Gaowei Zhang, Yi Wang.}
    {\textit{46th Annual Meeting of the Cognitive Science Society (CogSci)}, 2024.}
\end{enumerate}

\cvsection{HONORS AND AWARDS}{}
\begin{itemize}
  \awardentry{First-Class Scholarship}{Beijing University of Posts and Telecommunications}{2025}
  \awardentry{Deployed Application Award}{Innovative Applications of Artificial Intelligence (IAAI)}{2024}
  \awardentry{First Prize, Beijing Division}{China Undergraduate Mathematical Contest in Modeling}{2020}
  \awardentry{National Third Prize}{Chinese Mathematics Competitions}{2019}
\end{itemize}

\cvsection{SERVICES}{}
\begin{tabularx}{\textwidth}{@{}>{\bfseries\raggedright\arraybackslash}p{1.45in}@{\hspace{0.10in}}X@{}}
  Program Committee & AAAI 2026 AI for Social Impact Track \\
  Reviewer & NeurIPS 2026
\end{tabularx}

\cvsection{TEACHING EXPERIENCE}{}
\begin{itemize}
  \teachingentry{Social Interactions and Computing Systems}{Spring 2026}
  \teachingentry{Modern Software Development Methodologies and Practices}{Spring 2025}
  \teachingentry{Research Methods in Software Engineering}{Fall 2023}
  \teachingentry{Reinforcement Learning}{Spring 2023, Fall 2024}
\end{itemize}

\end{document}
